\documentclass[11pt]{article}
\usepackage[T1]{fontenc}
\usepackage[utf8]{inputenc}
\usepackage{lmodern}
\usepackage{amsmath,amsthm,amssymb}
\usepackage[margin=1in]{geometry}
\usepackage{microtype}
\usepackage[hidelinks]{hyperref}
\usepackage{url}
\hypersetup{
  pdftitle={A Lean Formalization of a Pairwise Common-Divisor Criterion for Ideals},
  pdfauthor={Lucas Freund},
  pdfsubject={Complete formalization of the polynomial-ideal criterion and its UFD form},
  pdfkeywords={Lean, Mathlib, prime avoidance, polynomial ideals}
}

\newtheorem{theorem}{Theorem}[section]
\newtheorem{lemma}[theorem]{Lemma}
\newtheorem{corollary}[theorem]{Corollary}
\theoremstyle{remark}
\newtheorem{remark}[theorem]{Remark}
\newcommand{\C}{\mathbb{C}}
\newcommand{\N}{\mathbb{N}}
\newcommand{\units}{R^{\times}}
\newcommand{\td}{\operatorname{tdeg}}
\newcommand{\pair}{\mathcal{P}}
\newcommand{\code}[1]{\texttt{\detokenize{#1}}}
\setlength{\emergencystretch}{2em}

\title{A Lean Formalization of a Pairwise\newline
Common-Divisor Criterion for Ideals}
\author{Lucas Freund\\[2pt]
  \small\texttt{lucasfreund2-cmd/open-math-agents}}
\date{27 November 2026\\[2pt]\small Formalization note}

\begin{document}
\maketitle

\begin{abstract}
Let $I$ be an ideal of a unique factorization domain. Suppose that every
pair of elements of $I$ admits a common divisor that is nonzero and not a
unit. Then a single prime element divides every element of $I$.
We give a Lean formalization of this classical argument and specialize
it to the polynomial-ideal question recorded as Mathematics Stack
Exchange question 2822876. The proof fixes one nonzero element of the
ideal, covers the ideal by the finitely many principal prime ideals
arising from its factors, and applies finite prime avoidance. An
additional lemma identifies positive total degree with being a nonzero
nonunit over a field. The development includes the zero ideal and zero
variables, uses the full pairwise hypothesis, and does not require a
Noetherian assumption. This note documents the mathematical argument,
its translation into Lean, and the reproducibility and review scope of
the accompanying artifact. The contribution is a formalization of
established mathematics.
\end{abstract}

\section{The source problem and its scope}

Schultz's 2018 question asks whether an ideal of
$\C[x_1,\ldots,x_n]$ whose elements have pairwise nontrivial common
divisors has a nonconstant divisor common to the whole ideal
\cite{schultz}. The same result appears as Exercise~5.5.3 in Appendix~A,
\emph{Basics from Algebraic Geometry}, of Kraft's notes, on printed
page~169 \cite{kraft}. In a 2021 MathOverflow answer, Sawin gives an
argument for unique factorization domains; a comment by \texttt{tj\_}
identifies finite prime avoidance as the underlying principle
\cite{sawin,tj}. These sources establish the prior mathematical context
of the present development. We make no claim of a new theorem or of
priority for its first formalization.

Write $A_n=\C[x_1,\ldots,x_n]$ and let $\td(f)$ denote the total degree,
with $\td(0)=0$. The precise polynomial statement is the following.

\begin{theorem}[Polynomial-ideal criterion]\label{thm:polynomial}
For every $n\in\N$ and ideal $I\subseteq A_n$, assume that
\begin{equation}\label{eq:polynomial-pair}
  \forall f,g\in I\quad
  \exists h\in A_n\quad
  \td(h)>0,\qquad h\mid f,\qquad h\mid g.
\end{equation}
Then there exists $h\in A_n$ such that
\begin{equation}\label{eq:polynomial-global}
  \td(h)>0
  \quad\text{and}\quad
  \forall f\in I,\ h\mid f.
\end{equation}
\end{theorem}

Both quantifiers in \eqref{eq:polynomial-pair} range over the entire
ideal, with no distinctness or nonzero restriction. The divisor may
depend on the pair; the conclusion removes that dependence. The
divisor is required to be a nonconstant polynomial, rather than merely
a nonunit, because the zero polynomial is a nonunit. Neither the
hypothesis nor the conclusion requires the divisor to belong to $I$.

The question's proposed route uses finite generation. The proof below
instead obtains finiteness from the factorization of one nonzero
element. This permits a statement for arbitrary ideals in an arbitrary
unique factorization domain (UFD). The specialization to complex
polynomials uses their UFD structure; algebraic closedness, analytic
properties of $\C$, and a chosen generating family play no role.

\section{A criterion in a unique factorization domain}

Throughout this section, $R$ is a commutative integral domain with
identity and unique factorization. A prime element is, by definition,
nonzero, not a unit, and prime for divisibility. We write $(p)=pR$.

\begin{theorem}[General criterion]\label{thm:ufd}
For an ideal $I$ of $R$, the following conditions are equivalent:
\begin{enumerate}
  \item Every $a,b\in I$ have a common divisor $d\ne0$ that is not a unit.
  \item There is a nonzero nonunit $d\in R$ with $d\mid a$ for every
    $a\in I$.
  \item There is a prime element $p\in R$ such that $I\subseteq(p)$.
\end{enumerate}
Thus the pairwise condition is equivalent to containment in a
\emph{nonzero principal prime ideal}.
\end{theorem}

The word ``nonzero'' in the last sentence is essential: the zero ideal
is itself a principal prime ideal in a domain, but its generator is
not a prime element. The theorem concerns containment, and does not
assert that $I$ itself is principal.

We recall the form of prime avoidance used in the proof. It is a
special case of the standard finite prime-avoidance lemma
\cite{stacks}; the short proof is included for completeness.

\begin{lemma}[Finite prime avoidance]\label{lem:avoidance}
Let $J,P_1,\ldots,P_r$ be ideals of a commutative ring, with every
$P_i$ prime. If $J\subseteq\bigcup_{i=1}^r P_i$, then
$J\subseteq P_i$ for some $i$.
\end{lemma}

\begin{proof}
Delete any ideal contained in another, leaving an incomparable family
with the same union. Suppose that $J$ is contained in none of them.
Choose $a_i\in J\setminus P_i$ and, for each $j\ne i$, choose
$b_{ij}\in P_j\setminus P_i$. Put
\[
  c_i=a_i\prod_{j\ne i}b_{ij},\qquad c=\sum_i c_i.
\]
Each $c_i$ lies in $J$, so $c\in J$. Primality gives $c_i\notin P_i$;
for $k\ne i$, the factor $b_{ki}$ gives $c_k\in P_i$. Consequently
$c\notin P_i$ for every $i$, contradicting the cover. Empty products
are $1$; an empty family cannot cover an ideal because every ideal
contains $0$.
\end{proof}

\begin{proof}[Proof of Theorem~\ref{thm:ufd}]
Assume condition~(1). If $I=(0)$, apply it to the pair $(0,0)$.
The resulting nonzero nonunit divides every element of $I$, proving
(2) in this case.

Otherwise choose $0\ne a\in I$. Condition~(1) applied to $(a,a)$
ensures that $a$ is not a unit. Let $p_1,\ldots,p_r$ represent its
distinct irreducible factors up to association. This is a finite
nonempty family, and each $p_i$ is prime because $R$ is a UFD.
For any $b\in I$, choose a common nonzero nonunit divisor $d$ of $a$
and $b$. An irreducible factor $q$ of $d$ divides both. Since $q\mid a$,
it is associated to one of the $p_i$, and hence $p_i\mid b$. We obtain
\begin{equation}\label{eq:cover}
  I\subseteq \bigcup_{i=1}^r(p_i).
\end{equation}
The ideals $(p_i)$ are prime. Lemma~\ref{lem:avoidance} gives
$I\subseteq(p_i)$ for some $i$, so (2) follows with $d=p_i$.

Next, assume (2). A nonzero nonunit $d$ has an irreducible factor
$p$, which is prime in a UFD. Divisibility is transitive, so
$p\mid a$ for every $a\in I$; this proves (3), including when $I=(0)$.
Finally, (3) implies (1) by using the same prime $p$ for every pair.
\end{proof}

The only finite set needed is the set of factors of $a$. No bound on
the number of generators of $I$ enters the proof. In
\eqref{eq:cover}, the index $i$ initially depends on $b$; ideal closure
and prime avoidance are precisely what allow one index to work for
all elements simultaneously.

\section{Nonconstant polynomials and boundary cases}

The passage from Theorem~\ref{thm:ufd} to
Theorem~\ref{thm:polynomial} requires the following elementary bridge.

\begin{lemma}\label{lem:degree}
For a field $K$, a variable type $\sigma$, and
$f\in K[X_s\mid s\in\sigma]$,
\[
  \td(f)>0
  \quad\Longleftrightarrow\quad
  f\ne0\ \text{and}\ f\text{ is not a unit}.
\]
\end{lemma}

\begin{proof}
A polynomial of positive total degree is nonzero, and units over a
field are nonzero constants, so it cannot be a unit. Conversely,
a polynomial of total degree zero is its constant coefficient.
If it is nonzero, that coefficient is an invertible field element;
the polynomial is therefore a unit. This contradicts the right-hand
side and proves the converse.
\end{proof}

Applying Lemma~\ref{lem:degree} to each divisor in
\eqref{eq:polynomial-pair}, followed by Theorem~\ref{thm:ufd}, yields
a global nonzero nonunit. A second application of the lemma gives
\eqref{eq:polynomial-global}. The same reasoning works over every
field. The submitted theorem retains $\C$ and finitely many variables
to match the source problem.

For $I=(0)$, extracting a divisor from the pair $(0,0)$ avoids any
unstated assumption that the ideal contains a nonzero element.
When $n\ge1$, a variable provides such a divisor. When $n=0$, the
polynomial ring is $\C$, so it has no positive-degree elements.
The pairwise hypothesis is then false, already at $(0,0)$, and the
implication remains valid. Thus admitting all $n\in\N$ does not create
a spurious existence claim about polynomials in zero variables.

Two further distinctions matter for the translation. First, checking
only pairs of generators is insufficient. In $K[x,y,z]$, the three
generators $xy,xz,yz$ have pairwise common nonconstant factors, but
they have none in common to all three. Their ideal also contains
$xy$ and $z(x+y)$, which have no common nonconstant factor. The source
hypothesis therefore fails for this ideal, as it should. Second, a
global divisor need not lie in the ideal: $x$ divides every element
of $(x^2,xy)\subset K[x,y]$, while $x\notin(x^2,xy)$, as every
polynomial in that ideal has zero component of total degree one.

Finally, having no common nonunit factor is a divisibility condition,
not a claim that two elements generate the unit ideal. For example,
$x$ and $y$ have no common nonconstant factor in $K[x,y]$, but
$(x,y)$ is proper. The formal statement uses divisibility directly,
avoiding an accidental strengthening to a B\'ezout condition.

\section{Lean architecture}

The companion source \cite{artifact} separates the general algebra
from the polynomial interpretation. The generic ring assumptions are
\code{CommRing R}, \code{IsDomain R}, and
\code{UniqueFactorizationMonoid R}. The four principal declarations
in namespace \code{OMA.MSE2822876} are:
\begin{enumerate}
  \item \code{common_divisor_of_pairwise}, proving the global-divisor
    conclusion in a UFD;
  \item \code{pairwise_iff_exists_prime}, proving equivalence with a
    prime element dividing the entire ideal;
  \item \code{positive_totalDegree_iff}, establishing
    Lemma~\ref{lem:degree} over any field;
  \item \code{polynomial_ideal_common_factor}, specializing the generic
    theorem to \code{MvPolynomial (Fin n)} with complex coefficients.
\end{enumerate}

For a nonzero element \code{a}, the proof enables Mathlib's strong
normalization structure and takes the finite set underlying
\code{normalizedFactors a}. Multiplicities are discarded because
only the available prime divisors matter. Mathlib's factorization
lemmas supply an irreducible factor of each common divisor and an
associated normalized factor of \code{a}. Association preserves
divisibility, so these representatives give exactly the cover in
\eqref{eq:cover} \cite{mathlib}.

The algebraic interface is
\code{Ideal.mem_span_singleton}: membership in the ideal generated
by $p$ is equivalent to divisibility by $p$.
\code{Ideal.isPrime_span_singleton_of_prime} supplies primality, and
\code{Ideal.subset_union_prime} supplies finite prime avoidance.
Within the \code{MvPolynomial} namespace, the degree bridge uses
\code{totalDegree_eq_zero_iff_eq_C} and
\code{isUnit_iff_totalDegree_of_isReduced}.
These library results are reused rather than reimplemented.

The proof uses classical reasoning to select factors and handle the
zero-ideal alternative. It establishes existence and supplies no
algorithm for computing a factor from an arbitrary presentation of an
ideal. This matches the existential source question. Polynomial
factorization algorithms, effective degree bounds, and extensions to
domains without unique factorization are outside this artifact's
claims.

\section{Verification, reproducibility, and review}

The artifact pins Lean~4.33.1 and the following dependencies:
\begin{center}
\small
\begin{tabular}{ll}
Mathlib & \texttt{0df444a360eaa60ab8c11dca51a86af692955474}\\[3pt]
Formal Conjectures & \texttt{9097e2f90e66d2f6244209cb1c07f7120d17ec53}
\end{tabular}
\end{center}
Formal Conjectures provides the pinned checking environment. The
question itself is newly stated in the companion
\path{statements/mse-2822876.lean}, rather than taken from an existing
upstream declaration. The proof is
\path{proofs/mse-2822876.lean}; its final wrapper is:
\begin{verbatim}
theorem OMA.solution :
    type_of% @OMA.S_mse_2822876.main := by
  intro n I hpair
  exact OMA.MSE2822876.polynomial_ideal_common_factor
    n I hpair
\end{verbatim}
Using the type of the statement fixes the target without using the
statement as a proved premise. The statement file retains the single
intentional \code{sorry} required by the site's translation-review
workflow. The solution file contains a complete proof. An axiom audit
of \code{OMA.solution} must confirm that the statement's proof hole
has not entered the derivation.

The unmodified repository checker accepted the proposed statement and
the complete, exact-target proof locally. For \code{OMA.solution},
Lean reported exactly \code{propext}, \code{Classical.choice}, and
\code{Quot.sound}, with no dependence on \code{sorryAx}. The verified
proof has 117 lines and contains neither \code{sorry} nor
\code{native_decide}. A separate AI-assisted audit checked the written
argument and the translation's boundary cases. The accompanying
\path{verification.json} records the source hashes, toolchain, checker
verdicts, and additional theorem audits.

For reproduction, use the pinned environment and the commands in the
artifact README, lint the statement and proof files, and run
\path{tools/check.sh} on the proof with the local statement available.
Retain the compiler output and resulting axiom report. Checking a
statement's elaboration alone does not validate its mathematical
translation, and checking a proof against a local statement does not
by itself establish acceptance by an external submission service.

At the preparation of this note, the statement-only submission is
upstream pull request~5285. Its GitHub compilation check accepted the
statement; human translation review remains pending \cite{submission}.
The complete proof is available separately in the author's repository
\cite{artifact}. This note makes no claim that LeftToSolve has marked
the problem solved or merged the proof. The artifact provides a
complete formal argument; review of fidelity to the source remains a
separate part of the submission process.

\section*{Authorship and AI assistance}

This note and the accompanying formalization were prepared for Lucas
Freund's repository with AI assistance from ChatGPT, GPT-6 Astra Pro.
The assistance covered source investigation, proof development,
mathematical checking, and drafting. The Lean kernel checks the formal
derivation; it does not establish literary originality, validate
source attribution, or confer peer review. The work is presented as a
formalization note with an inspectable artifact, without a claim of
journal publication or peer-reviewed acceptance.

\begin{thebibliography}{9}
\small
\raggedright

\bibitem{schultz}
L.~Schultz,
\emph{Is there a non-constant polynomial that divides every element of
this given ideal?}
Mathematics Stack Exchange, question~2822876, 17 June 2018.
\url{https://math.stackexchange.com/questions/2822876/}.
Question text licensed CC BY-SA~4.0; the statement above is a
mathematical restatement.

\bibitem{kraft}
H.~Kraft,
\emph{Basics from Algebraic Geometry}, Appendix~A, Exercise~5.5.3,
printed p.~169. Online notes; date of this PDF version not established.
\href{https://dmi.unibas.ch/fileadmin/user_upload/dmi/Personen/Kraft_Hanspeter/Algebraic_Geometry.pdf}{University of Basel PDF}.

\bibitem{sawin}
W.~Sawin,
Answer to MathOverflow question~410693, 13 December 2021.
\url{https://mathoverflow.net/a/410694}.

\bibitem{tj}
\texttt{tj\_},
Comment on Sawin's answer identifying the prime-avoidance argument,
13 December 2021.
\url{https://mathoverflow.net/questions/410693/}.

\bibitem{stacks}
The Stacks Project Authors,
\emph{The Stacks Project}, Lemma~10.15.2, Prime avoidance, Tag~00DS.
\url{https://stacks.math.columbia.edu/tag/00DS}.

\bibitem{mathlib}
The Mathlib Community,
\emph{Mathlib}, commit
\texttt{0df444a360eaa60ab8c11dca51a86af692955474}.
\href{https://github.com/leanprover-community/mathlib4/tree/0df444a360eaa60ab8c11dca51a86af692955474}{Pinned source tree}.

\bibitem{artifact}
L.~Freund,
\emph{Complete formalization of MSE 2822876}, accompanying Lean source,
repository \texttt{lucasfreund2-cmd/open-math-agents},
commit \texttt{f13994c69223125c88eddaf72aafc21786a369c5}.
\href{https://github.com/lucasfreund2-cmd/open-math-agents/blob/f13994c69223125c88eddaf72aafc21786a369c5/proofs/mse-2822876.lean}{Immutable proof source}.

\bibitem{submission}
L.~Freund,
Statement submission for \texttt{mse-2822876},
\emph{open-math-agents}, pull request~5285, 27 November 2026.
\url{https://github.com/GuillaumeBAUD-coder/open-math-agents/pull/5285}.

\end{thebibliography}
\end{document}
